% !TEX root = ../PhD.tex
%
% Bridge section: not part of the original publication.
% Describes the RDF-Connect pipelines behind MuMo, as found in
% mumo-full/pipeline/pipeline/{pipeline.ttl,auth-pipeline.ttl} with
% @rdfc/sds-processors-ts 2.0.2-alpha.1.

\section{Discussion in the context of this dissertation}\label{sec:mumo-dissertation}

\noindent\textit{%
    This section is not part of the original publication.
    It describes the RDF-Connect pipelines that MuMo runs on,
    which the publication leaves implicit,
    and relates the deployment to the research questions of this dissertation.
}

\subsection{Contribution to the research questions}

The publication describes MuMo from the point of view of the museums that use it.
Its data path (Section~\ref{sec:implementation}) starts with ``a measurement is ingested'',
and leaves out how a measurement becomes part of a Linked Data Event Stream.
That step is an RDF-Connect pipeline (Chapter~\ref{ch:rdf-connect}).
MuMo is therefore where this dissertation tests RQ1 and RQ2 outside a demonstration:
the pipeline runs unattended, polls its sources every five seconds,
and keeps its state across restarts.
MuMo does not test RQ3, since it shows nothing about how the pipeline description was written.

\subsection{The MuMo pipelines}

The data pipeline has two branches, one for sensors and one for observations,
which run as eleven processor instances in a single Node.js runner.
Each branch has the same five stages.
\begin{enumerate}
    \item \texttt{MumoFetch} polls the history endpoint of the MuMo dashboard
        and saves its position, so that a restart resumes instead of reprocessing.
    \item \texttt{MumoMapper} maps the dashboard's JSON to RDF:
        platforms and sensors in SOSA with their current location in CIDOC-CRM,
        and observations following ISO~19156 with QUDT units (Section~\ref{sec:datamodel}).
    \item \texttt{Sdsify} wraps every member in an \texttt{sds:Record}
        for a stream identifier named in the pipeline description,
        together with a SHACL shape that delimits each member.
    \item \texttt{Bucketize} fragments the stream according to strategies declared in RDF.
        Observations are fragmented by the location of the sensor's platform, by platform,
        and by time (year, month, day, hour, and minute);
        sensors are fragmented by location.
        This is the fragmentation hierarchy of Section~\ref{fragmentation-as-a-design-lever},
        and therefore also the unit of access control
        (Section~\ref{policy-aligned-fragmentation-ldes-solid-together}).
    \item \texttt{LdesDiskWriter} writes the fragments to disk as an LDES.
\end{enumerate}
The governance layer is a second RDF-Connect pipeline.
It starts the Community Solid Server as a processor that serves the directory the data pipeline writes to,
the ACL generator of Section~\ref{runtime-components} as another,
and a processor that issues WebIDs.

For RQ1, the deployment shows how little of a pipeline is specific to its project.
Of the processor types in the data pipeline,
only \texttt{MumoFetch} and \texttt{MumoMapper} were written for MuMo:
they talk to the dashboard and map its JSON.
\texttt{Sdsify}, \texttt{Bucketize}, \texttt{LdesDiskWriter}, and \texttt{GlobRead}
are taken unchanged from the processor library,
the same processors that the pipelines of Section~\ref{sec:usecases} use.
Everything that concerns streaming, fragmentation, and publication is configuration.
In particular, the fragmentation that the publication treats as a central design lever
is a list of strategies in the pipeline description:
changing how data is fragmented, and with it what can be shared,
changes RDF and not code.

\subsection{Provenance in the deployment}

For RQ2, the deployment attaches provenance at two levels, through SDS (Section~\ref{sec:sds}).
Every published item carries a record that names its stream and the buckets it belongs to,
because \texttt{Bucketize} moves each record from its incoming stream to the stream it produces.
Every published stream carries a description of how it was derived:
\texttt{Bucketize} reads the description of its incoming stream,
adds a \texttt{p-plan:Activity} of type \texttt{ldes:Bucketization}
that used the incoming stream and the fragmentation strategies, together with its start time,
and declares the new stream \texttt{prov:wasGeneratedBy} that activity.
A consumer of a MuMo LDES can therefore find out,
from the data alone, which stream an observation belongs to and how that stream was fragmented.

The deployment does not yet reach the full chain of Section~\ref{sec:rdfc-sds}, for three reasons.
First, the description of each incoming stream is read from a static file
instead of being produced by \texttt{Sdsify}, so the first link is written by hand.
Second, the stream identifiers in the deployed configuration do not line up:
the identifier \texttt{Sdsify} assigns is described nowhere,
and the bucketizer's output stream reuses the identifier of its input,
so the derivation refers to itself.
Third, the activity in a stream description is a fresh blank node
rather than the processor in the orchestrator's PROV-O trace,
and the deployment does not write that trace.
Each of these is a matter of configuration or a small change to the SDS processors,
and none requires a change to the framework.
\todo{If these three issues are fixed before submission, replace this paragraph with the resulting chain from observation to pipeline.}
